\documentclass[11pt]{article}
\usepackage{amsmath,amssymb,booktabs}
\usepackage[margin=1in]{geometry}

\title{Nominal Velocity Model of the Unitree Go2}
\date{October 2026}

\begin{document}
\maketitle

\section{General Model}
Each command segment $k$ holds a forward speed $v_x \ge 0$ and a yaw rate $\omega$ for a duration $T$.
The yaw command is split into its left and right parts, $\omega^{+} = \max(\omega, 0)$ and $\omega^{-} = \min(\omega, 0)$.
The body-frame velocity the robot achieves is modelled as
\begin{equation}
  \begin{bmatrix} \tilde v_x \\ \tilde\omega^{+} \\ \tilde\omega^{-} \end{bmatrix}
  =
  \underbrace{\begin{bmatrix} g_x & 0 & 0 \\ 0 & g_\omega^{+} & 0 \\ 0 & 0 & g_\omega^{-} \end{bmatrix}}_{\mathbf{G}_3}
  \begin{bmatrix} v_x \\ \omega^{+} \\ \omega^{-} \end{bmatrix}
  + \frac{1}{T}
  \begin{bmatrix} o_x\,\mathbf{1}(v_x>0) \\ o_\omega^{+}\,\mathbf{1}(\omega>0) \\ o_\omega^{-}\,\mathbf{1}(\omega<0) \end{bmatrix}
  +
  \begin{bmatrix} 0 \\ \kappa\, v_x \\ 0 \end{bmatrix},
  \label{eq:general}
\end{equation}
and the yaw rate of the robot is $\tilde\omega = \tilde\omega^{+} + \tilde\omega^{-}$.
Here $g_x$ and $o_x$ are the forward gain and offset, $g_\omega^{\pm}$ and $o_\omega^{\pm}$ the left ($+$) and right ($-$) turn gains and offsets, $\mathbf{1}(\cdot)$ equals 1 when its condition holds and 0 otherwise, and $\kappa$ is the yaw drift per metre of commanded forward walking.

\section{Turning from Standstill and after Running}
The robot needs time to begin a turn, and this time depends on whether its gait is already running.
We therefore distinguish the case $c_k$ of every in-place turn segment ($v_x = 0$, $\omega \neq 0$):
\begin{equation}
  c_k =
  \begin{cases}
    \mathrm{st} & \text{first segment of a turn, robot at standstill } (k = 0 \text{ or } \mathbf{u}_{k-1} = \mathbf{0}),\\
    \mathrm{run} & \text{first segment of a turn, right after running } (\mathbf{u}_{k-1} \neq \mathbf{0}),\\
    \mathrm{cont} & \text{later segment of the same turn}.
  \end{cases}
\end{equation}
The turn gains and offsets in \eqref{eq:general} take the values of that case:
\begin{equation}
  g_\omega^{\pm} = g_{\omega,c_k}^{\pm}, \qquad o_\omega^{\pm} = o_{\omega,c_k}^{\pm},
  \qquad c_k \in \{\mathrm{st}, \mathrm{run}, \mathrm{cont}\}.
\end{equation}
A turn that starts from standstill loses about $\lambda_{\mathrm{st}} = 0.30$\,s of its first segment and one that starts right after running about $\lambda_{\mathrm{run}} = 0.10$\,s.
Losing this time at the steady turn rate is equivalent to a smaller gain in the first segment,
\begin{equation}
  g_{\omega,c}^{\pm} = g_{\omega,\mathrm{cont}}^{\pm}\left(1 - \frac{\lambda_c}{T}\right),
  \qquad c \in \{\mathrm{st}, \mathrm{run}\}.
\end{equation}

\section{Parameter Values}
\begin{table}[h]
  \centering
  \begin{tabular}{lll}
    \toprule
    Variable & Value & Description \\
    \midrule
    $g_x$ & 1.05 & forward gain \\
    $o_x$ & 0 m & forward offset per segment \\
    $\kappa$ & 0.05 rad/m & walking yaw drift \\
    $g_{\omega,\mathrm{st}}^{+}$, $g_{\omega,\mathrm{st}}^{-}$ & 0.49, 0.49 & turn gains, first segment from standstill \\
    $o_{\omega,\mathrm{st}}^{+}$, $o_{\omega,\mathrm{st}}^{-}$ & 0 rad, 0 rad & turn offsets, first segment from standstill \\
    $g_{\omega,\mathrm{run}}^{+}$, $g_{\omega,\mathrm{run}}^{-}$ & 0.63, 0.63 & turn gains, first segment after running \\
    $o_{\omega,\mathrm{run}}^{+}$, $o_{\omega,\mathrm{run}}^{-}$ & 0 rad, 0 rad & turn offsets, first segment after running \\
    $g_{\omega,\mathrm{cont}}^{+}$, $g_{\omega,\mathrm{cont}}^{-}$ & 0.70, 0.70 & turn gains, later segments of a turn \\
    $o_{\omega,\mathrm{cont}}^{+}$, $o_{\omega,\mathrm{cont}}^{-}$ & 0 rad, 0 rad & turn offsets, later segments of a turn \\
    $T$ & 1 s & segment duration \\
    \bottomrule
  \end{tabular}
  \caption{Model parameters identified from 30 pilot runs recorded on 2026-10-14.}
\end{table}

\end{document}
